\begin{definition}[Unary-admissible partial hidden sources]
\label{def:typed-unary-admissible-source}
A finite unary-functionally CLOSED PARTIAL terminal hidden
structure \(C\) is \emph{unary-admissible} if, for EVERY
\(x\in C\), the FULL induced hidden diagram on
\(\cl_C\{x\}\), including all named total unary functions,
their literal term equalities and the binary relations
internal to this closure, is isomorphic to an ACTUAL
completed one-generated induced hidden member.
This does NOT assume \(C\) itself is complete or
homomorphism-embedding completable.

The condition is substantive. Arbitrary finite partial
structures with total unary symbols may have arbitrarily
large one-generated closures and violate the named
quotient-transition identities; mere unary closedness
does NOT imply unary-admissibility.
\end{definition}

\begin{example}[Arbitrarily long malformed unary transition cycles]
\label{ex:typed-malformed-unary-long-cycles}
The distinction in
\cref{def:typed-unary-admissible-source}
already occurs in the ACTUAL two-original-sort
common-block SAP+CCO structure
\cref{ex:typed-mixed-block-orbital-diagonal}.
Its proper quotient sorts $Q_A,Q_B$ carry the
NAMED inverse transitions
$f=\pi_{r,E_A}:Q_A\to Q_B$
and $g=\pi_{r^\smile,E_B}:Q_B\to Q_A$,
with $g f=\operatorname{id}_{Q_A}$
in every COMPLETED hidden member.

For arbitrary $n\ge2$, take a finite
PARTIAL total-unary hidden source
with quotient vertices $x_0,\ldots,x_{n-1}$
in $Q_A$ and $y_0,\ldots,y_{n-1}$
in $Q_B$, no real vertices and
NO displayed binary orbitals.
Put $f(x_i)=y_i$ and
$g(y_i)=x_{i+1\pmod n}$.
The generated transition orbit of ANY vertex
is the WHOLE $2n$-vertex structure,
so it has NO proper nonempty
unary-closed induced substructure and
is IRREDUCIBLE. It has no hidden
homomorphism-embedding completion:
function preservation would imply
$h(x_{i+1})=g f(h(x_i))=h(x_i)$,
contradicting injectivity on this irreducible source.
Thus no uniform bound on NONCOMPLETABLE
unary-CLOSED substructures can hold for
ARBITRARY total-unary partial inputs.

These malformed sources fail
HN Definition~2.8(ii): their
one-generated IRREDUCIBLE closures
are NOT completed hidden members.
In the actual HN test, (ii) enforces the
unary-admissibility needed for the
RESTRICTED bounded-obstruction proof;
the terminal Ramsey claim is NOT
refuted by this example.
\end{example}

\begin{lemma}[Uniform typed unary closure in genuine hidden diagrams]
\label{lem:typed-unary-closure-bound}
There is a constant \(m_{\mathrm{typ}}\), depending only on
the finite terminal hidden language, such that the unary
functional closure of one vertex has at most
\(m_{\mathrm{typ}}\) vertices in EVERY ACTUAL completed
hidden member. Consequently the same bound holds for
every one-generated closure in a
\cref{def:typed-unary-admissible-source} source.
NO uniform bound is asserted for arbitrary total-unary
PARTIAL source structures.
\end{lemma}

\begin{proof}
A unary coordinate term is a finite composition of projections, upward
quotient maps and quotient transitions.  Suppose a transition induced by
\(R\) carries an \(E\)-coordinate to the transported
\(E^R\)-coordinate and is followed by the upward map to some represented
\(H\geq E^R\).  The lattice isomorphism in
\cref{lem:typed-quotient-transport} has a unique pullback
\(H_0\geq E\) with \(H_0^R=H\); meet-irreducibility is preserved by
transport, so this pullback is again a represented coordinate.  On quotient
classes
\[
  u_{E^R,H}\circ\pi_{R,E}
   =\pi_{R,H_0}\circ u_{E,H_0}.
\]
Thus every upward map occurring after a transition can be commuted one step
towards the source.  Repeating this operation collects all upward maps at the
source.  Consequently every coordinate term is a single projection to a
represented quotient (if the source is real), or a single initial upward map
(if it is already a coordinate), followed by a transition-only composite.
The explicit composition and
upward-interchange identities
\cref{lem:typed-transition-composition-normal-form}
show that every
transition-only composite
is the quotient map of
one actual relation
in a finite \(\Rel^{ab}\)
family, or an existing
identity term if that
RAW relational
representative contains
logical identity.
The latter case uses
\cref{lem:typed-identity-transport-redundant},
not an artificial STRONG
identity-labelled atom.
There are only finitely
many represented coordinate sorts.  Thus only finitely many unary term
functions can occur from any fixed hidden sort, uniformly in the finite
structure, which gives the asserted bound.
\end{proof}

The next lemma records the more precise normal form which will be used in the
amalgamation argument.

\begin{definition}[Transition-closed root signature]
\label{def:typed-transition-root-signature}
Let two real roots \(x\in a\), \(y\in b\) carry a guarded pair constraint
\(D\in\Dat^{ab}\).  Its canonical root signature contains, for every
represented \(J\in\mathcal I_a\) with
\(DD^{\smile}\subseteq J\), the direct quotient-transition identity
\[
 \pi_{D,J}(F_J(x))=F_{J^D}(y),
\]
together with its full functional closure.  If the constraint is represented
by a theta bundle
\(D=\bigcap_iP_i\) with \(P_i\in\Rch^{ab}\), we additionally include the
corresponding transition identities forced by each branch profile \(P_i\).
Thus both intersection-forced and branch-forced quotient data are retained.
\end{definition}

The signature is compatible with every actual root pair.  More generally, if
non-empty definable pair relations \(R,S\in\Rel^{ab}\) both contain an
atomic type \(c\), and both are quotient-functional from the same represented
\(J\), then \(J^R=J^c=J^S\) and the three quotient maps agree on every
source \(J\)-class.  Indeed an actual \(c\)-neighbour of any representative
computes the unique source-class label for all three relations.  Hence the
direct exact-domain transition and all branch transitions in one canonical
theta root signature are consistent.

\begin{lemma}[Typed coordinate-term normal form]
\label{lem:typed-coordinate-normal-form}
Let \(\tau\) be a unary term whose source is a real sort \(a\) and whose
value is either the source real point itself or a hidden quotient coordinate.
Then there is a definable equivalence \(K_\tau\) on \(a\), with
\(K_\tau\) equality exactly for the identity real-point term and otherwise a
proper meet-irreducible, such that for all real \(x,x'\in a\),
\[
 \tau(x)=\tau(x')\quad\Longleftrightarrow\quad xK_\tau x'.
\tag{QT-NF}\label{eq:typed-coordinate-normal-form}
\]
Moreover the value \(\tau(x)\) is obtained from the
\(K_\tau\)-coordinate of \(x\) by a finite chain of quotient transitions and upward quotient maps.
Consequently the stabiliser of \(\tau(x)\) in any automorphism group equals
the stabiliser of the \(K_\tau\)-class of \(x\).
\end{lemma}

\begin{proof}
The identity real-point term has kernel equality.  A non-identity
real-source term first applies a named projection
\(F_{E_0}:a\to B_{a,E_0}\) to a represented proper
meet-irreducible \(E_0\).  The rest is a finite word in named
upward quotient maps and quotient transitions (including their
named inverses).

We prove the stronger LITERAL normal form, including its named maps.
Whenever an upward map \(u_{F,H}\) follows a quotient transition
\(\pi_{R,E}\), the source-sort pullback \(H_0\ge E\) is again a
represented proper meet-irreducible, by the interval isomorphism
\cref{lem:typed-quotient-transport}.  Replace these two
consecutive maps, in source-to-target order, by
\[
 \pi_{R,E}\ \text{then}\ u_{F,H}
 \quad=\quad
 u_{E,H_0}\ \text{then}\ \pi_{R,H_0},
\]
as literal unary functions by the upward-interchange identity
\eqref{eq:typed-upward-interchange}.  Each replacement reduces
the finite number of transition-before-upward inversions in
the word.  Once no such inversion remains, combine the
upward maps into the initial projection \(F_K\), where
\(K\) is a represented proper meet-irreducible on the
ORIGINAL source sort \(a\).  Combine all remaining
transitions by the actual RAW composition identity
\eqref{eq:typed-transition-composition}.  If the resulting
RAW relation contains logical identity, the corresponding
quotient map is the already named identity by
\cref{lem:typed-identity-transport-redundant};
otherwise it is the named transition \(\pi_{T,K}\) for
an ACTUAL \(T\in\Rel^{ab}\).  Write \(\theta\) for
the named transition in the second case, and for the
EXISTING named identity in the first.  Thus we have
a literal equality of term functions
\[
                       \tau=\theta\circ F_K.
\]
Every transition is a quotient-set bijection, so
\(\ker(\tau)=K\), giving \eqref{eq:typed-coordinate-normal-form}.
The value of \(\tau(x)\) is thus obtained by a NAMED
term from the \(K_\tau\)-coordinate, even if intermediate
transitions change original sorts.

Finally the quotient transition is equivariant and bijective.
An automorphism fixes \(\tau(x)\) exactly when it fixes the
\(K_\tau\)-class of \(x\), proving the stated stabiliser equality.
\end{proof}

\begin{corollary}[Named factorisation of forced typed real-source terms]
\label{cor:typed-forced-term-factorisation}
Let \(J\) be a represented proper meet-irreducible on original
sort \(a\), and let \(\tau\) be an ACTUAL unary term from \(a\)
with kernel \(K_\tau\).  If \(J\le K_\tau\), then there is a
literal NAMED unary term
\(\sigma_{\tau,J}:B_{a,J}\to\operatorname{cod}(\tau)\) such that
\[
                \tau(x)=\sigma_{\tau,J}(F_J(x))
                \qquad(x\in a).
\]
No choice of groupoid chart, quotient representative or
original-sort ordering is involved.
\end{corollary}
\begin{proof}
Since \(J>\Id_a\), the real identity term cannot satisfy
\(J\le K_\tau\).  For a proper hidden-coordinate term,
the preceding literal normal form gives
\(\tau=\theta\circ F_{K_\tau}\),
where \(\theta\) is a NAMED quotient transition or the
existing identity term.  The upward map
\(u_{J,K_\tau}\) is named, and
\(F_{K_\tau}=u_{J,K_\tau}\circ F_J\).
Thus \(\sigma_{\tau,J}=\theta\circ u_{J,K_\tau}\)
is a literal term in the ACTUAL hidden language,
also for RAW diagonal-return composites.
\end{proof}

\begin{lemma}[Typed return transport]\label{lem:typed-return-transport}
Let \(P_1,\ldots,P_n\) be a composable typed word from a sort \(a_0\)
back to \(a_0\), and suppose
\[
 P_1\cdots P_n\subseteq E_0
\]
for a definable equivalence \(E_0\) on \(a_0\). Put
\(R_i=P_1\cdots P_i\). Every prefix transports \(E_0\)
to an invariant equivalence \(E_i\), every factor induces one
quotient bijection, and the loop has IDENTITY holonomy.
For EVERY represented proper meet-irreducible
\(J\supseteq E_0\), the corresponding intermediate
kernels are NAMED proper represented quotient coordinates
and the actual quotient return is the identity on \(a_0/J\).
\end{lemma}
\begin{proof}
Apply the general ORIGINAL invariant-return theorem
\cref{lem:foundational-typed-return-holonomy}
to the actual equality-inclusive RAW typed relations \(P_i\).
It canonically gives
\[
 E_i=R_i^\smile E_0R_i,\qquad
 J_i=R_i^\smile JR_i,\qquad
 q_n\circ\cdots\circ q_1=\operatorname{id}_{a_0/J}.
\]
The transported \(J_i\) are meet-irreducible in the FULL
invariant interval by
\cref{lem:typed-quotient-transport}.
In this terminal typed source EVERY invariant equivalence is
FULL RAW exact by
\cref{lem:typed-terminal-all-equivalences-exact}.
Also, visible SAP and
\cref{lem:no-real-algebraic-transition}
prevent a nonidentity represented proper coordinate from
transporting to OLD real equality; the interval isomorphism
preserves strictness below the universal equivalence.
Thus the \(J_i\) are precisely the NAMED represented proper
quotient coordinates of the terminal hidden signature,
and each quotient bijection is its existing named unary
transition (with actual identity returns interpreted as
the named identity term). The canonical quotients in the
foundational construction use the SAME \(J_i\) at both
macro sides incident at the intermediate OLD root.
No transport chart or completion-dependent seam equality
has been chosen.
\end{proof}

The transition closure corrects the first typed failure mechanism as well.
If an atomic type \(t:a\to b\) and a represented meet-irreducible \(J\) on
\(a\) satisfy \(tt^{\smile}\subseteq J\), then a \(b\)-point \(y\)
uniquely determines the \(J\)-class of every \(t\)-neighbour: project \(y\)
to its proper \(J^t\)-coordinate and apply the inverse of \(\pi_{t,J}\).
Properness of \(J^t\) follows from
\cref{lem:no-real-algebraic-transition}.

\begin{remark}[The mixed-block diagnostic]
The two-sorted mixed-block structure is the basic test which invalidates the
naive ``carry sort indices through'' proof.  Let \(E_A,E_B\) be the
same-block equivalences on the two real sorts and let \(r:A\to B\) be the
same-block cross type.  Then
\[
 rr^{\smile}=E_A,
 \qquad
 E_A^r=E_B,
\]
so the hidden transition \(\pi_{r,E_A}\) sends the represented
\(E_A\)-class of an \(A\)-point to the represented \(E_B\)-class of its
\(r\)-neighbour.  The return word \(rr^{\smile}\) therefore transports the
quotient from \(A/E_A\) to \(B/E_B\) and back with identity holonomy; no
cross-sort factor is asserted to lie in a same-sort equivalence.  Thus both
failure mechanisms of the naive typed proof are visible, and repaired, in
this smallest test case.
\end{remark}

